% 数学ノート集・共通テンプレート v2（積分論2026年度第1部の版面を参考）
% A5判。章内で定義・定理・例・注意等の番号を共通化。
% コンパイル： platex -interaction=nonstopmode math_note_template_v2.tex (2回)
%              dvipdfmx math_note_template_v2.dvi
\documentclass[a5paper,10pt,onecolumn]{jsarticle}
\usepackage[dvipdfmx,top=22mm,bottom=20mm,inner=19mm,outer=18mm,headheight=13pt,headsep=6mm]{geometry}
\usepackage{amsmath,amssymb,amsthm,mathtools}
\usepackage{enumitem,fancyhdr,titlesec}
\usepackage[dvipdfmx,bookmarks=true,colorlinks=false]{hyperref}
\usepackage{pxjahyper}

% ----- 差し替える文書情報 -----
\newcommand{\NoteTitle}{多様体上の微分形式とStokesの定理}
\newcommand{\NoteSeries}{数学ノート集}

% ----- 文字・組版 -----
% 参考PDFの質素な明朝＋Computer Modernの印象を引き継ぎながら，
% A5判で窮屈にならないよう本文の行間を少し確保する．
\linespread{1.13}
\setlength{\parindent}{1zw}
\setlength{\parskip}{0pt}
\setlength{\abovedisplayskip}{6pt plus 2pt minus 2pt}
\setlength{\belowdisplayskip}{6pt plus 2pt minus 2pt}
\setlength{\abovedisplayshortskip}{3pt plus 2pt minus 1pt}
\setlength{\belowdisplayshortskip}{4pt plus 2pt minus 1pt}
\setlist[enumerate]{topsep=3pt,itemsep=1pt,parsep=0pt,leftmargin=2.4zw}
\titleformat{\section}{\normalfont\normalsize\bfseries}{\thesection}{0.8zw}{}
\titlespacing*{\section}{0pt}{2.0ex plus .4ex}{.9ex}
\titleformat{\subsection}{\normalfont\normalsize\bfseries}{\thesubsection}{0.7zw}{}
\titlespacing*{\subsection}{0pt}{1.5ex plus .3ex}{.6ex}
\raggedbottom

% ----- 柱とページ番号：ページ上部、罫線なし -----
\fancypagestyle{mathnote}{%
  \fancyhf{}%
  \fancyhead[L]{\footnotesize \NoteSeries\ （\NoteTitle）}%
  \fancyhead[R]{\footnotesize \thepage}%
  \renewcommand{\headrulewidth}{0pt}%
  \renewcommand{\footrulewidth}{0pt}%
}
\pagestyle{mathnote}

% ----- 節内通し番号の定義・定理群 -----
\newtheoremstyle{mathnotetheorem}
 {4.5pt}{4.5pt}{\normalfont}{}%
 {\bfseries}{}{0.65zw}%
 {\thmname{#1}\thmnumber{ #2}\thmnote{（#3）}}
\theoremstyle{mathnotetheorem}
\newtheorem{theorem}{定理}[section]
\newtheorem{definition}[theorem]{定義}
\newtheorem{lemma}[theorem]{補題}
\newtheorem{proposition}[theorem]{命題}
\newtheorem{corollary}[theorem]{系}
\newtheorem{example}[theorem]{例}
\newtheorem{remark}[theorem]{注意}
\renewcommand{\proofname}{証明}
\renewcommand{\qedsymbol}{$\square$}
\makeatletter
\renewenvironment{proof}[1][\proofname]{%
  \par\pushQED{\qed}\normalfont\topsep4\p@\@plus2\p@\relax
  \trivlist\item[\hskip\labelsep\bfseries #1]\ignorespaces
}{\popQED\endtrivlist\@endpefalse}
\makeatother

% ----- 使うノートで適宜拡張する記号 -----
\newcommand{\R}{\mathbb{R}}
\newcommand{\norm}[1]{\lVert #1\rVert}
\newcommand{\abs}[1]{\lvert #1\rvert}

\begin{document}
\section*{読了後にできること}
本稿を読むための前提は，微分形式，向き，Stokesの定理．
本稿の具体的な到達課題は，Stokesの定理を具体的な境界積分で検算する．

\section{余接空間と微分形式}

本稿では滑らかな多様体の接空間・接束，
Euclid空間上の外微分と引き戻し，
多変数積分の変数変換公式を前提とする．
多様体はHausdorffかつ第二可算とし，
境界を持つ場合も含める．

\begin{definition}
$p\in M$ に対し $T_p^*M$ を
接空間 $T_pM$ の実線形双対空間という．
$k$-次微分形式とは，$p\in M$ ごとに
$\omega_p\in\Lambda^kT_p^*M$ を対応させ，
局所座標 $(x^1,\dots,x^n)$ で
\[
 \omega=\sum_{i_1<\cdots<i_k}
 a_{i_1\cdots i_k}(x)\,
 dx^{i_1}\wedge\cdots\wedge dx^{i_k}
\]
と書いたとき全係数が滑らかであるものをいう．
全体を $\Omega^k(M)$ と書く．
\end{definition}

\begin{proposition}[外微分の座標独立性]
\label{prop:d-global}
各座標近傍で
\[
 d\left(\sum_Ia_I\,dx^I\right)
 =\sum_{I,j}\frac{\partial a_I}{\partial x^j}
 dx^j\wedge dx^I
\]
と定義すると，局所定義は座標変換によって一致し，
大域的な線形写像
$d:\Omega^k(M)\to\Omega^{k+1}(M)$ を定める．
また
$d^2=0$ および
\[
 d(\omega\wedge\eta)
 =d\omega\wedge\eta+(-1)^k\omega\wedge d\eta
\]
が成り立つ．
\end{proposition}
\begin{proof}
座標変換 $\phi$ に対して
$d(\phi^*\omega)=\phi^*(d\omega)$ が
Euclid空間の微分形式で成立することを用いる．
この等式は連鎖律と二階偏導関数の交換法則を
座標表示に代入して直接確認できる．
従って重なり合う座標近傍では
外微分の表示が同じ形式を定めるため，
局所定義を貼り合わせて大域的な $d$ が得られる．
残る二つの等式はEuclid空間上で
係数の微分と交代性を用いて成立するので，
全ての座標近傍上で成立し，
従って多様体上でも成立する．
\end{proof}

\section{引き戻しと変数変換}

\begin{definition}
滑らかな写像 $F:M\to N$ と
$\omega\in\Omega^k(N)$ に対して
\[
 (F^*\omega)_p(v_1,\dots,v_k)
 =\omega_{F(p)}(dF_pv_1,\dots,dF_pv_k)
\]
を引き戻しという．
\end{definition}

\begin{proposition}[引き戻しの函手性]\label{prop:pullback}
滑らかな写像 $F:M\to N$，$G:N\to P$ について
\[
 (G\circ F)^*=F^*\circ G^*,\qquad
 d(F^*\omega)=F^*(d\omega)
\]
が成り立つ．
また $F^*(\omega\wedge\eta)
=F^*\omega\wedge F^*\eta$ である．
\end{proposition}
\begin{proof}
第一式は
$d(GF)_p=dG_{F(p)}dF_p$ を
引き戻しの定義式へ代入して得る．
外積については，交代多重線形形式の外積を定める
有限個の置換の和へ $dF_p$ を各引数に適用すればよい．
外微分との可換性は局所座標に還元すると
Euclid空間上の同じ公式になるので，
命題\ref{prop:d-global}の議論から成立する．
\end{proof}

\section{向き付けと体積形式}

\begin{definition}
$n$ 次元多様体 $M$ のAtlasが
任意の二座標の遷移写像について
Jacobi行列式が正となるとき，
これを\emph{向き付けられたAtlas}という．
そのようなAtlasを持つ多様体を向き付け可能という．
\end{definition}

\begin{lemma}[向き付けと正の体積形式]
\label{lem:volume}
滑らかな $n$ 次元多様体 $M$ が向き付け可能なら，
各点で零でない滑らかな $n$-形式 $\nu$ が存在する．
逆に各点で零でない $n$-形式は向き付けを定める．
\end{lemma}
\begin{proof}
正の遷移行列式を持つ向き付けAtlasを取り，
それに従属する単位の分割 $\rho_i$ を選ぶ．
各座標近傍で
$\omega_i=dx_i^1\wedge\cdots\wedge dx_i^n$
を取り，
\[
 \nu=\sum_i\rho_i\omega_i
\]
と置く．各 $\rho_i\omega_i$ は台の外で零と延長し，
局所有限性により $\nu$ は滑らかである．
向き付けられた座標で見ると，
各 $\omega_i$ は正のJacobi行列式を通じた
正の倍数になるので，
$\rho_i\ge0$ かつ $\sum_i\rho_i=1$ から
$\nu$ は各点で正の非零形式となる．
逆に非零な $n$-形式 $\nu$ があれば，
$\nu$ が局所座標の標準体積形式の正の倍数となる
座標だけを選ぶ．それらの遷移行列式は正だから
向き付けAtlasを得る．
\end{proof}

\section{微分形式の積分}

\begin{definition}
$M$ を向き付けられた $n$ 次元多様体とする．
$\omega$ をコンパクト台を持つ $n$-形式とする．
向き付けられた座標近傍に従属する
単位の分割 $\{\rho_i\}$ を取る．
$\rho_i\omega$ の座標表示を
$f_i(x)\,dx^1\wedge\cdots\wedge dx^n$ としたとき
\[
 \int_M\omega=\sum_i\int_{\varphi_i(U_i)}f_i(x)\,dx
\]
と定める．
\end{definition}

\begin{proposition}[積分の定義の独立性]
\label{prop:integral-independent}
上の積分は，向き付けられた座標近傍と
それに従属する単位の分割の取り方によらない．
\end{proposition}
\begin{proof}
二つの座標近傍の重なりで，
$n$-形式の座標係数はJacobi行列式を掛けて変換される．
向き付けられた座標変換では行列式が正なので，
Lebesgue積分の変数変換公式から
\[
 \int_{\varphi(U\cap V)}(\varphi^{-1})^*\omega
 =\int_{\psi(U\cap V)}(\psi^{-1})^*\omega
\]
が成り立つ．
二組の単位の分割 $\{\rho_i\}$，$\{\eta_j\}$ を取ると
$\sum_{i,j}\rho_i\eta_j=1$ であり，
コンパクト台のため有限個の項だけが積分に寄与する．
各交わりの座標で上の変数変換公式を適用すれば，
二組の分割で定義した積分が共に
$\sum_{i,j}\int_M\rho_i\eta_j\omega$ に等しくなる．
従って定義は選択によらない．
\end{proof}

\section{境界付き多様体と誘導向き}

\begin{definition}
$n$ 次元\emph{境界付き滑らかな多様体}とは，
局所座標が閉半空間
\[
 \mathbb H^n=\{x\in\mathbb R^n:x_n\ge0\}
\]
の相対開集合への同相写像であり，
座標変換が半空間の近傍へ滑らかに延長できるものをいう．
$x_n=0$ に対応する点全体を境界 $\partial M$ と書く．
\end{definition}

\begin{definition}
$M$ を向き付けられた境界付き $n$ 次元多様体とする．
$\partial M$ の向きは，境界における外向き接ベクトル
$\nu_{\mathrm{out}}$ と
$v_1,\dots,v_{n-1}\in T_p(\partial M)$ に対し，
\[
 (\nu_{\mathrm{out}},v_1,\dots,v_{n-1})
\]
が $T_pM$ の正の基底となるとき，
$(v_1,\dots,v_{n-1})$ を正とする向きによって定める．
これを外向き法線優先の向きという．
\end{definition}

\begin{proposition}\label{prop:boundary-orient}
$\partial M$ は自然に $(n-1)$ 次元滑らかな多様体となる．
$M$ が向き付けられていれば上の定義は
境界の向きをwell-definedに定める．
\end{proposition}
\begin{proof}
境界付き多様体の座標変換は境界を境界へ写すので，
$x_n=0$ の部分へ制限した座標変換は
$\mathbb R^{n-1}$ の開集合間の滑らかな微分同相となる．
従って境界の座標近傍は
$(n-1)$ 次元滑らかなAtlasをなす．
外向き接ベクトルは境界に接するベクトルに
正の外向き成分を足したものを取り，
その違いは境界接ベクトルの加算と正のスカラー倍だけである．
行列式の符号はこれらの変更で変化しないので，
誘導向きは外向きベクトルの選択によらない．
座標変換でも向き付けは保たれるため，
境界の向きはwell-definedである．
\end{proof}

\section{Euclid空間でのStokes公式}

\begin{lemma}[全空間での微分形式の積分]\label{lem:stokes-full}
$\omega\in\Omega_c^{n-1}(\mathbb R^n)$ なら
\[
 \int_{\mathbb R^n}d\omega=0.
\]
\end{lemma}
\begin{proof}
$\omega$ を
\[
 \omega=\sum_{i=1}^n(-1)^{i-1}
 a_i(x)\,dx^1\wedge\cdots
 \wedge\widehat{dx^i}\wedge\cdots\wedge dx^n
\]
と書く．
帽子はその因子を除くことを意味する．
外微分の定義から
\[
 d\omega=
 \left(\sum_{i=1}^n\partial_i a_i\right)
 dx^1\wedge\cdots\wedge dx^n.
\]
各 $a_i$ はコンパクト台を持つので，
Fubiniの定理と微積分の基本定理から
\[
 \int_{\mathbb R^n}\partial_i a_i(x)\,dx
 =\int_{\mathbb R^{n-1}}
 \left(\int_{-\infty}^\infty
   \partial_i a_i(x)\,dx_i\right)dx_{\widehat i}=0.
\]
有限和を取って結論を得る．
\end{proof}

\begin{lemma}[半空間でのStokes公式]\label{lem:stokes-half}
$\mathbb H^n=\{x_n\ge0\}$ に標準向きを与える．
コンパクト台を持つ滑らかな $(n-1)$-形式
$\omega$ に対して
\[
 \int_{\mathbb H^n}d\omega
 =\int_{\partial\mathbb H^n}\omega
\]
が成り立つ．右辺の境界には外向き法線優先の向きを入れる．
\end{lemma}
\begin{proof}
補題\ref{lem:stokes-full}と同様に
\[
 \omega=\sum_{i=1}^n(-1)^{i-1}
 a_i\,dx^1\wedge\cdots
 \wedge\widehat{dx^i}\wedge\cdots\wedge dx^n
\]
と書く．
$i<n$ の項を $x_i$ について積分すると，
コンパクト台とFubiniの定理によって零になる．
$i=n$ の項だけは
\[
 \int_0^\infty\partial_n a_n(x',t)\,dt=-a_n(x',0)
\]
なので
\[
 \int_{\mathbb H^n}d\omega
 =-\int_{\mathbb R^{n-1}}a_n(x',0)\,dx'.
\]
一方境界へ制限した $\omega$ は
\[
 \omega|_{x_n=0}
 =(-1)^{n-1}a_n(x',0)
 dx^1\wedge\cdots\wedge dx^{n-1}
\]
である．
外向き法線は $-\partial_n$ であり，
外向き法線優先の境界向きは，
標準的な $(x_1,\dots,x_{n-1})$ の向きの
$(-1)^n$ 倍となる．
従って
\[
 \int_{\partial\mathbb H^n}\omega
 =(-1)^n(-1)^{n-1}
 \int_{\mathbb R^{n-1}}a_n(x',0)\,dx'
 =-\int a_n(x',0)\,dx',
\]
左辺と一致する．
\end{proof}

\section{一般Stokesの定理}

\begin{theorem}[Stokesの定理]\label{thm:stokes}
$M$ を向き付けられた $n$ 次元境界付き滑らかな多様体とし，
$\omega\in\Omega_c^{n-1}(M)$ とする．
このとき
\[
 \boxed{\displaystyle
 \int_Md\omega=\int_{\partial M}\omega}
\]
が成立する．
\end{theorem}
\begin{proof}
$M$ の向き付けられた座標近傍からなる開被覆を取り，
それに従属する滑らかな単位の分割
$\{\rho_i\}$ を選ぶ．
$\omega$ はコンパクト台なので，
積分に寄与する項は有限個であり，
$\omega=\sum_i\rho_i\omega$ と書ける．
外微分は線形だから
$d\omega=\sum_i d(\rho_i\omega)$ である．
従って各 $\rho_i\omega$ について公式を示せば十分である．

$\rho_i\omega$ は一つの座標近傍内にコンパクト台を持つ．
その近傍が境界に接しない場合には，
Euclid空間の開集合に座標変換できる．
コンパクト台なので零延長を行い，
補題\ref{lem:stokes-full}から
$\int_Md(\rho_i\omega)=0$ を得る．
境界上では $\rho_i\omega=0$ なので
右辺も零である．

座標近傍が境界に接する場合には，
半空間 $\mathbb H^n$ の相対開集合に座標変換する．
コンパクト台のため零延長した滑らかな
$(n-1)$-形式を半空間全体に考えられる．
座標変換が外微分と可換であることと，
向き付けられた変数変換公式を用い，
補題\ref{lem:stokes-half}を適用すると
\[
 \int_Md(\rho_i\omega)
 =\int_{\partial M}\rho_i\omega
\]
となる．
これを有限個の $i$ について足し合わせれば
\[
 \int_Md\omega
 =\sum_i\int_Md(\rho_i\omega)
 =\sum_i\int_{\partial M}\rho_i\omega
 =\int_{\partial M}\omega
\]
が得られる．
\end{proof}

\section{古典的な積分定理との関係}

\begin{proposition}[Greenの公式]\label{prop:green}
$D\subset\mathbb R^2$ をコンパクトな滑らかな境界を持つ領域，
$P,Q$ をその近傍の滑らかな関数とする．
$\partial D$ に正の境界向きを入れると
\[
 \oint_{\partial D}P\,dx+Q\,dy
 =\iint_D\left(\frac{\partial Q}{\partial x}
               -\frac{\partial P}{\partial y}\right)dx\,dy.
\]
\end{proposition}
\begin{proof}
$\omega=P\,dx+Q\,dy$ と置くと
\[
 d\omega=(Q_x-P_y)\,dx\wedge dy
\]
である．
定理\ref{thm:stokes}を二次元領域 $D$ に適用すると
表示式を得る．
\end{proof}

\begin{proposition}[発散定理]\label{prop:divergence}
$D\subset\mathbb R^3$ を滑らかなコンパクト領域，
$F=(F_1,F_2,F_3)$ を滑らかなベクトル場とする．
外向き単位法線を $n$ とすると
\[
 \int_{\partial D}F\cdot n\,dS
 =\int_D\operatorname{div}F\,dV.
\]
\end{proposition}
\begin{proof}
\[
 \omega=F_1\,dy\wedge dz+
 F_2\,dz\wedge dx+
 F_3\,dx\wedge dy
\]
と置く．
外微分を計算すると
$d\omega=(\partial_xF_1+\partial_yF_2+\partial_zF_3)
 dx\wedge dy\wedge dz$ である．
境界へ $\omega$ を引き戻したものは，
向き付けられた接ベクトル $u,v$ に対して
$F\cdot(u\times v)$ に等しいので，
その積分は外向き法線への流束積分
$\int_{\partial D}F\cdot n\,dS$ に一致する．
Stokesの定理を適用すれば結論を得る．
\end{proof}

\begin{proposition}[古典的Stokes公式]\label{prop:curl}
向き付けられた滑らかな曲面 $S\subset\mathbb R^3$ と
その境界 $\partial S$ に対して，
滑らかなベクトル場 $F$ が
\[
 \oint_{\partial S}F\cdot d\boldsymbol r
 =\int_S(\operatorname{curl}F)\cdot n\,dS
\]
を満たす．
\end{proposition}
\begin{proof}
$\omega=F_1dx+F_2dy+F_3dz$ と置く．
微分計算により
\[
 d\omega=(\partial_yF_3-\partial_zF_2)\,dy\wedge dz
 +(\partial_zF_1-\partial_xF_3)\,dz\wedge dx
 +(\partial_xF_2-\partial_yF_1)\,dx\wedge dy.
\]
これは曲面に引き戻すと
$(\operatorname{curl}F)\cdot n\,dS$ に一致する．
境界への $\omega$ の引き戻しは
線積分の被積分形式である．
定理\ref{thm:stokes}を二次元多様体 $S$ に適用すれば
表示式を得る．
\end{proof}

\begin{remark}
一般Stokesの定理は境界付き多様体上の外微分と積分の関係であり，
Green・Gauss・古典Stokesの公式を同じ等式の
異なる座標表現として統一する．
\end{remark}

\section{到達課題の完全解答}
\noindent\textbf{円板上のStokesの定理．}
\begin{proposition}
$D=\{(x,y):x^2+y^2\le1\}$ に標準の向きを入れ，
$\alpha=x\,dy$ とする．このとき
\[
\int_{\partial D}\alpha
=\int_Dd\alpha=\pi
\]
である．
\end{proposition}
\begin{proof}
$d\alpha=dx\wedge dy$ なので
$\int_Dd\alpha$ は単位円板の面積 $\pi$ である．
境界を反時計回りに
$\gamma(t)=(\cos t,\sin t)$
（$0\le t\le2\pi$）と置くと
\[
\int_{\partial D}\alpha
=\int_0^{2\pi}\cos t\cdot\cos t\,dt=\pi.
\]
双方が一致するのでStokesの定理を検算できる．
\end{proof}
\noindent\textbf{到達確認．}
$\beta=-y\,dx$ の境界積分を求めよ．
\par\noindent\textbf{解答．}
$d\beta=-dy\wedge dx=dx\wedge dy$ だから
$\int_{\partial D}\beta=\pi$ である．
\end{document}